\section*{Capítulo \ref{ch:b3:colloids} --- Interfaces, tensoativos e coloides}

\begin{solution}{exo:b3:colloids:1}
$r = \qty{0.50}{\micro m}$: $\Delta p = 2 \times 0.0720/\num{5.0e-7} = \qty{2.9e5}{Pa}$; no interior, $1.0 + 2.9 = \qty{3.9}{bar}$.
\end{solution}

\begin{solution}{exo:b3:colloids:2}
$\ln(p/p^*) = 2 \times 0.0720 \times \num{1.807e-5}/(\num{1.0e-8} \times 8.314 \times 298.15) = 0.105$; $p/p^* = 1.11$.
\end{solution}

\begin{solution}{exo:b3:colloids:3}
$\cos\theta = (40 - 20)/72 = 0.278$, $\theta = \ang{74}$: abaixo de $90^\circ$, o líquido \omterm{def:b3:colloids:wetting}{molha}
parcialmente o sólido.
\end{solution}

\begin{solution}{exo:b3:colloids:4}
NaCl: $I = \qty{10}{mol.m^{-3}}$, $\kappa^{-1} = 0.96\sqrt{100/10} = \qty{3.0}{nm}$. \ce{MgSO4}: $I = \qty{40}{mol.m^{-3}}$, \qty{1.5}{nm}.
\end{solution}

\begin{solution}{exo:b3:colloids:5}
$\Gamma = \num{12.6e-3}/(2 \times 8.314 \times 298.15) = \qty{2.54e-6}{mol.m^{-2}}$; área $1/(\Gamma N_A) =
\qty{0.65}{nm^2}$ por molécula.
\end{solution}

\begin{solution}{exo:b3:colloids:6}
Abaixo da CMC, a inclinação é de $-12.9$ \unit{mN/m} por unidade de $\log c$ (de $-5.0$ a $-4.0$); o
patamar fica em 33.4. A reta o atinge em $\log c = -4.0 + (39.1 - 33.4)/12.9 = -3.56$: CMC
$\approx \qty{2.8e-4}{mol/L}$.
\end{solution}

\begin{solution}{exo:b3:colloids:7}
SDS: $P = 0.350/(0.60 \times 1.67) = 0.35$, na fronteira entre esferas e cilindros curtos:
\omterm{def:b3:colloids:micelle}{micelas} esféricas perto da CMC. Lipídio: $P = 0.700/(0.65 \times 1.67) = 0.64$:
bicamadas, logo vesículas.
\end{solution}

\begin{solution}{exo:b3:colloids:8}
\ce{CaCl2}: $60/64 = \qty{0.94}{mM}$ de \ce{Ca^{2+}}; \ce{AlCl3}: $60/729 = \qty{0.082}{mM}$ de \ce{Al^{3+}}. O \ce{Na3PO4}
coagula por meio de seu \ce{Na+} (a cerca de \qty{20}{mM} de sal); o fosfato, um coíon,
pode até se adsorver e aumentar a carga negativa.
\end{solution}

\begin{solution}{exo:b3:colloids:9}
Parte hidrofílica $6 \times 44.0 + 17.0 = \qty{281.0}{g/mol}$ de \qty{450.0}{g/mol} (massas atômicas do livro):
HLB $= 20 \times 281.0/450.0 = 12.5$, um emulsificante óleo em água.
\end{solution}

\begin{solution}{exo:b3:colloids:10}
O óleo das gotículas pequenas é mais solúvel na fase contínua (pelo fator
$\exp(2\gamma V_{\mathrm m}/rRT)$, com um expoente dez vezes maior para \qty{50}{nm} do que para \qty{500}{nm}): ele
se difunde para as gotículas grandes, que crescem enquanto as pequenas desaparecem.
Esse \omterm{def:b3:colloids:ripening}{amadurecimento de Ostwald} engrossa uma \omterm{def:b3:colloids:colloid}{emulsão} mesmo que duas gotículas nunca se fundam.
\end{solution}

\begin{solution}{exo:b3:colloids:11}
O menisco é uma calota esférica de raio $r/\cos\theta$: o líquido logo abaixo dele fica a uma
pressão menor de $2\gamma\cos\theta/r$, e a coluna sobe até que $\rho gh$ compense essa diferença. Para a água,
$h = 2 \times 0.072/(997 \times 9.81 \times \num{1.0e-4}) = \qty{0.15}{m}$.
\end{solution}

\begin{solution}{exo:b3:colloids:12}
A \qty{100}{mM}, o \omterm{def:b3:electrode-kinetics:double-layer}{comprimento de Debye} é inferior a \qty{1}{nm} e a repulsão se extingue antes
da atração: não resta barreira (além da curva de \qty{50}{mM}). Um polímero livre não age
na superfície; uma camada de polímero enxertado repele no contato, qualquer que seja o
sal, pela entropia perdida quando as cadeias se sobrepõem.
\end{solution}

\begin{solution}{pb:b3:colloids:1}
\textbf{1.} Substituições na rede cristalina (\ce{Al^{3+}} no lugar de \ce{Si^{4+}}, \ce{Mg^{2+}} no lugar de \ce{Al^{3+}}) deixam
uma carga negativa permanente, compensada por cátions trocáveis.
\textbf{2.} $I = \frac12(1.0 + 1.0) = \qty{1.0}{mM}$.
\textbf{3.} $\kappa^{-1} = \qty{9.6}{nm}$.
\textbf{4.} $a\kappa = 10$: a dupla camada é fina em comparação com a partícula, a
condição da aproximação de Derjaguin.
\textbf{5.} $v = 2 \times 1650 \times 9.81 \times (\num{1.0e-7})^2/(9 \times \num{0.89e-3}) = \qty{4.0e-8}{m/s}$, cerca de \qty{3.5}{mm} por dia.
\textbf{6.} Suas \omterm{def:b3:electrode-kinetics:double-layer}{camadas difusas} se repelem antes que a atração de van der Waals prevaleça.
\textbf{7.} $V = 2\pi\varepsilon_r\varepsilon_0a\psi^2\ln(1 + \eu^{-\kappa h}) - Aa/12h$.
\textbf{8.} Cerca de $39\,kT$.
\textbf{9.} Uma fração da ordem de $\eu^{-39} \approx 10^{-17}$ das colisões: praticamente nunca.
\textbf{10.} Ela desaparece: toda colisão leva ao contato.
\textbf{11.} Um \omterm{def:b3:electrode-kinetics:double-layer}{comprimento de Debye} menor confina a repulsão a distâncias pequenas, onde
a atração, que cresce como $1/h$, domina.
\textbf{12.} A CCC varia como $z^{-6}$ da carga do contraíon.
\textbf{13.} $50/64 = \qty{0.78}{mM}$.
\textbf{14.} $50/729 = \qty{0.069}{mM}$.
\textbf{15.} Os cátions são os contraíons das partículas negativas: eles se acumulam na
dupla camada e a blindam.
\textbf{16.} O sulfato é um coíon, repelido pelas partículas.
\textbf{17.} \qty{0.069}{mol} de \ce{Al^{3+}} por metro cúbico.
\textbf{18.} \qty{0.034}{mol} de $\ce{Al2(SO4)3}$.
\textbf{19.} $0.0343 \times 342.3 = \qty{11.7}{g}$.
\textbf{20.} \qty{11.7}{kg} por hora.
\textbf{21.} $\ce{Al^3+ + 3 HCO3- -> Al(OH)3 + 3 CO2}$.
\textbf{22.} $3 \times 0.069 = \qty{0.21}{mM}$ do \qty{1.0}{mM}: um quinto. O hidrogenocarbonato
restante e o \ce{CO2} dissolvido tamponam a água: o pH cai apenas
moderadamente.
\textbf{23.} As espécies de alumínio muito carregadas se adsorvem e invertem a carga das
partículas, que voltam a se repelir.
\textbf{24.} \textbf{Cerca de \qty{12}{g} de sulfato de alumínio por metro cúbico} (\qty{11.7}{g}), antes do
excesso necessário, na prática, para a hidrólise.
\end{solution}
